%------------------------------------------------------------------------------%
\documentclass[a4paper,12pt,fleqn]{article}

\usepackage{calculo_varias_variaveis-1}

\ShowAnswers

%------------------------------------------------------------------------------%
\begin{document}

\makeHeader{CFVVI}{Prova 5 -- Turma 7}{19/12/2025}

% 01
%------------------------------------------------------------------------------%
\exercise{25\%}
Considere a equação em coordenadas polares
\(
  \left(\frac{r}{6}\right)^2 = \left(4\cos^2\theta+9\sin^2\theta\right)^{-1},
\)
escreva-a em coordenadas cartesianas, identifique a cônica e esboce seu gráfico.

\begin{questiononly}
  \bigskip
  \hfill\includegraphics{axis_xy}
\end{questiononly}

\clearpagequestiononly

\begin{answer}
\begin{multicols}{2}
\begin{align*}
  \left(\frac{r}{6}\right)^2        & = \left(4\cos^2\theta+9\sin^2\theta\right)^{-1} \\[2mm]
  \frac{r^2}{36}                    & = \frac{1}{4\cos^2\theta+9\sin^2\theta}         \\[2mm]
  r^2                               & = \frac{36}{4\cos^2\theta+9\sin^2\theta}
\end{align*}
\vspace{-1.5\baselineskip}
\begin{align*}
  r^2\left(4\cos^2\theta+9\sin^2\theta\right) & = 36 \\
  4r^2\cos^2\theta + 9r^2\sin^2\theta         & = 36 \\
  4x^2 + 9y^2                                 & = 36 \\
  \frac{4x^2}{36} + \frac{9y^2}{36}           & = 1  \\
  \frac{x^2}{3^2} + \frac{y^2}{2^2}           & = 1
\end{align*}

\columnbreak
A cônica é uma elípse
\bigskip

\includegraphics{T7_A_conica}
\end{multicols}

\end{answer}

% 02
%------------------------------------------------------------------------------%
\exercise{25\%}
Determine todas as raízes cúbicas do número complexo
\(
    z = -4\left(1+i \sqrt{3}\right)
\)
\clearpagequestiononly

\begin{answer}
O módulo de $z$ é
\[
    \abs{z}
    = \sqrt{(-4)^2 + \left(-4\sqrt{3}\right)^2}
    = \sqrt{16 + 16\cdot 3}
    = \sqrt{64}
    = 8
\]
Derminamos o argumento \(\varphi = \arg(z)\) resolvendo o sistema
\begin{align*}
    \sin(\varphi) & = \frac{\Im(z)}{\abs{z}} = \frac{-4\sqrt{3}}{8} = \frac{-\sqrt{3}}{2} \\
    \cos(\varphi) & = \frac{\Re(z)}{\abs{z}} = \frac{-4}{8} = \frac{-1}{2}
\end{align*}
portanto
\[
    \varphi = \arg(z) = \left(2n - \frac{2}{3}\right)\pi \qquad n\in\Z
\]
Podemos escolher
\[
  \varphi = -\frac{2\pi}{3}
  \qquad\text{ ou }\qquad
  \varphi = \frac{4\pi}{3}
\]
Vou seguir com
\(
  \varphi = -\frac{2\pi}{3}
\)

Pela segunda fórmula de Moivre, as raízes cúbicas de $z$ são dadas por
\begin{align*}
    u_k
    & = \sqrt[3]{8}\left(\cos\frac{\varphi+2k\pi}{3}
                      + i\sin\frac{\varphi+2k\pi}{3}\right) & & \quad k = 0, 1, 2 \\
    & = 2 \left(\cos\frac{\sfrac{-2\pi}{3}+2k\pi}{3}
             + i\sin\frac{\sfrac{-2\pi}{3}+2k\pi}{3}\right)\\
    & = 2 \left(\cos\frac{-2\pi+6k\pi}{9}
             + i\sin\frac{-2\pi+6k\pi}{9}\right)
\end{align*}
\[
    u_0=2\left(\cos\frac{-2\pi}{9}+i\sin\frac{-2\pi}{9}\right)
\]
\[
    u_1=2\left(\cos\frac{4\pi}{9}+i\sin\frac{4\pi}{9}\right)
\]
\[
    u_2=2\left(\cos\frac{10\pi}{9}+i\sin\frac{10\pi}{9}\right)
\]
\end{answer}


% 03
%------------------------------------------------------------------------------%
\exercise{50\%}
Determine os pontos candidatos a extremos de
\(
  x^2+e^{y}
\)
sujeita à restrição
\(
  x^2+y=1
\)
\clearpagequestiononly

\begin{answer}
\begin{multicols}{2}
Queremos os extremos de
\[
  f(x, y) = x^2 + e^{y}
\]
sujeita à restrição
\[
  g(x, y) = x^2 + y - 1 = 0
\]
Gradientes
\begin{align*}
  \nabla f & = \vetor{2x}{e^{y}} \\
  \nabla g & = \vetor{2x}{1}
\end{align*}

O sistema de Lagrange é
\begin{align*}
  2x    & =2\lambda x \\
  e^{y} & =\lambda    \\
  x^2+y & =1
\end{align*}

Da primeira equação
\[
  x = 0 \quad \text{ou} \quad \lambda=1
\]

Caso $x=0$, da restrição
\[
  y = 1
\]
e da segunda equação
\[
  \lambda = e
\]
obtermos o primeiro ponto
\[
  P_1 = (0, 1)
\]

Caso $\lambda=1$, da segunda equação
\[
  e^{y}=1 \quad\Rightarrow\quad y=0
\]
Da restrição
\[
x^2=1 \quad\Rightarrow\quad x=\pm1
\]
Encontramos o segundo e terceiro pontos
\[
P_2 = (1, 0) \qquad P_3 = (-1, 0)
\]
\end{multicols}
\end{answer}

%------------------------------------------------------------------------------%
\end{document}
%------------------------------------------------------------------------------%


